\documentclass[11pt,a4paper]{article}

\usepackage[T1]{fontenc}
\usepackage[utf8]{inputenc}
\usepackage{lmodern}
\usepackage[margin=1in]{geometry}
\usepackage{booktabs}
\usepackage{amsmath}
\usepackage{graphicx}
\usepackage{pgfplots}
\usepackage[numbers]{natbib}
\usepackage[hidelinks]{hyperref}
\usepackage{url}

\pgfplotsset{compat=1.18}

% One visual language for every figure: no vertical rules, no chartjunk, a faint
% horizontal grid, and the same three-colour ramp everywhere so a model keeps its colour
% from one plot to the next.
\definecolor{anvilA}{HTML}{4C6EB1}
\definecolor{anvilB}{HTML}{C9772E}
\definecolor{anvilC}{HTML}{3F7D5A}
\definecolor{anvilD}{HTML}{8A5A9E}
\definecolor{anvilE}{HTML}{9E3B4E}

\pgfplotsset{
  anvilbar/.style={
    ybar,
    width=\linewidth,
    height=6.2cm,
    bar width=3.6pt,
    ymin=0, ymax=1.05,
    ytick={0,0.25,0.5,0.75,1.0},
    ymajorgrids=true,
    grid style={gray!25},
    axis lines*=left,
    tick align=outside,
    enlarge x limits=0.10,
    legend style={draw=none, fill=none, font=\footnotesize, at={(0.5,1.06)},
                  anchor=south, legend columns=-1, /tikz/every even column/.append style={column sep=4pt}},
    every axis plot/.append style={draw=none},
    x tick label style={font=\small},
    y tick label style={font=\small},
  }
}

\title{Anvil: executable verification of LLM-generated\\HPC operational artifacts}
% No institutional affiliation: the work was done independently, with no supervision or
% review from anyone at a university, and naming one would put the reading of an
% endorsement on a party that never gave it.
\author{Antonio Rotundo\\
  \small Independent researcher}
% Fixed, not `\today`: on arXiv the date would be whichever day their build ran, and
% the PDF here is built byte-reproducibly under SOURCE_DATE_EPOCH, so a floating date
% was the one thing that made two builds of the same source differ.
\date{August 15, 2026}

\begin{document}
\maketitle

\begin{abstract}
Assistants that write SLURM job scripts for supercomputer users are evaluated with similarity
metrics, because no benchmark verifies whether the artifacts they produce actually work. Anvil
verifies them by execution against a declared reference cluster: five levels covering syntax,
submittability, execution, effective resource fit and safety, bracketed by an oracle that must
score $1.0$ and a deliberately faulty model that must score $0.0$. On eight from-scratch tasks and
$44$ induced repair tasks we measure five models across three families and two generations of one
of them. Submittability is not ordered by model size: on the from-scratch set a 3B model of one
family and a 12B of another lead it while a 7B trails both, and within one family the level falls
as size rises. Real job submission looked nearly redundant against a \texttt{bash} sandbox, one
verdict in $3120$ artifacts, until a fifth model wrote \texttt{srun} inside its scripts: the
sandbox has no allocation for a job step to attach to, and the count moved to six in $3900$ with
$32$ samples failing in the sandbox and passing on the scheduler. On a task set whose memory need
only the payload knows, the same comparison reads $288$ of $900$. Retrieved documentation costs
the smaller model $28$ points of resource fit whatever the
text says, while the larger one is unaffected by irrelevant text and loses $6$ points to the one
document that addresses the level it concerns. A task built so that GNU coreutils and
\texttt{uutils} must judge it differently is failed by all three models asked for reasons that
precede the difference. We also report, rather than conceal, three measurement faults found and
corrected
during the work, one of which had already been published.
\end{abstract}

\section{Introduction}

A wrong \verb|--mem| does not look wrong. It looks plausible, and the job dies at hour six.

Operational HPC artifacts, batch scripts and container recipes, sit in a gap. Execution-based
evaluation is standard for program synthesis and increasingly common for infrastructure code, yet
the assistants appearing for HPC user support are still scored by how closely their output
resembles a reference answer. Resemblance is the wrong instrument here for a specific reason: the
failures that matter are invisible in the text. A directive placed after the first command is
silently ignored by the scheduler. An omitted \verb|--ntasks| is not an absent request but a
request for a different number. A memory figure one factor too small parses, submits, and runs
until it does not.

Unlike cloud infrastructure code, where a plan cannot be executed cheaply and evaluation falls back
on graph edit distance or model judges, these artifacts are directly checkable.
\verb|sbatch --test-only| costs milliseconds, the job really runs, and the ground truth is
objective. The empty slot is not a hard problem; it is an unoccupied one.

This paper reports what the resulting instrument says about five models and, with more care, what
it says about itself, including one claim that four models supported and a fifth overturned.

\section{Related work}

The works below were selected from a sweep of fourteen queries against OpenAlex, Crossref, DBLP
and arXiv, run by \texttt{scripts/litsweep.py} and returning $761$ unique records. Its full output
is committed with this manuscript, so the selection can be audited rather than taken on trust.

\paragraph{Assistants for HPC users, and how they are judged.} Support assistants for
supercomputer users are being built and deployed: a chatbot for autonomous user support
\citep{pasupuleti2025chatbot}, a science-gateway assistant whose outcomes were reported from a
national laboratory \citep{biggs2024gateway}, and a Slurm-native serving infrastructure for running
such models on the cluster itself \citep{doosthosseini2024chatai}. \citet{zheng2026comparative}
compare several of these approaches. The task predates the models: the NERSC job script generator
\citep{shao2022nersc} produced batch scripts from a form, with correctness guaranteed by
construction rather than measured. What is missing across this line is an execution-grounded way to
say whether the artifacts these systems emit would run, which is the gap this work addresses.

\paragraph{Language models for HPC code.} A parallel line targets the source code rather than the
operational artifacts around it: modelling parallel programs \citep{nichols2023hpccoder}, zero-shot
parallelization guided by graph neural networks \citep{mahmud2023autoparllm}, and MPI code
generation for HPC workloads \citep{valerolara2026chatmpi}. These evaluate a kernel that compiles
and runs. A batch script is judged by a scheduler instead, and fails in ways a compiler has no
opinion about: a directive after the first command is accepted and ignored, and an omitted one is
not an absent request but a different one.

\paragraph{Execution-grounded evaluation elsewhere.} Outside HPC, execution is the norm.
\citet{chen2021codex} established the unbiased $\mathrm{pass}@k$ estimator this work uses
unchanged. Infrastructure-as-code is the closest neighbour and shows the constraint: unable to
execute a plan cheaply, \citet{santos2026terraform} benchmark generated Terraform through static
analysis. HPC operational artifacts do not share that constraint. \texttt{sbatch -{}-test-only}
costs milliseconds and the job really runs, which is why the ground truth here can be objective
where theirs cannot.

\section{The benchmark}

\paragraph{T1, writing from scratch.} Eight tasks, each a support request in ordinary words paired
with a machine-readable specification: node and task counts, walltime and memory bounds, GPUs,
directives the task demands explicitly, and strings the script must print when it runs.

\paragraph{T2, diagnosing and repairing.} $44$ tasks derived mechanically from the T1 canonical
solutions by seven fault injectors, anchored to classes observed on a real model rather than
invented: an omitted directive whose SLURM default hides it, a misplaced directive, prose leaking
into a value, an absent no-default directive, a script stripped of every directive, a payload that
no longer matches its own specification, and a value the scheduler rejects. Induction keeps only
the variants that actually fail verification, since an injector producing an accidentally valid
script is a bug in the injector and not a fault worth teaching a model to repair. Repairs are
graded by the same verifier as T1, with no partial credit for being closer to correct.

\section{Method}

\paragraph{Five levels.} The strict metric requires all of them.

\begin{description}\setlength{\itemsep}{2pt}
  \item[syntax] a shebang, \texttt{bash -n}, and no directive after the first command.
  \item[submittability] \texttt{sbatch -{}-test-only} accepts it.
  \item[functional] it runs and exits zero.
  \item[resource fit] the \emph{effective} request matches the specification, with SLURM's
        documented defaults applied.
  \item[safety] destructive-pattern probes, which gate execution.
\end{description}

\paragraph{A skipped level is never a passed one.} Without a scheduler, submittability is skipped
and scored as not passed. One refinement was forced by measurement: a level skipped because
\emph{this machine} lacks a capability is pass-through, since every artifact is skipped alike,
whereas a level skipped because \emph{this artifact} cannot be judged anywhere is a failure. Before
the distinction existed, ten repairs of a dependency task passed strict by being unjudgeable.

\paragraph{Effective requests, not string presence.} A serial script omitting \verb|--nodes| still
requests one node and is correct. An early version checked for the string and failed scripts the
scheduler accepts, which is precisely the surface-form matching this work exists to replace.

\paragraph{Both bounds shipped.} An oracle returning canonical solutions must score $1.0$ or the
benchmark is defective; a broken model returning faulty artifacts must score $0.0$ strict, and its
sampling is arranged so the destructive flavour is always drawn, because a guard that never fires
is decoration.

\paragraph{A declared reference cluster.} Four nodes, sixteen cores and one GPU each, in a
container running a live \texttt{slurmctld}, so that scores do not depend on the hardware of
whoever runs the benchmark. Two preflights protect it: a canary that distinguishes a failing
scheduler from a failing model, and a topology canary that refuses to score submittability when the
scheduler in front of it is not the declared cluster.

\paragraph{Generation and verification are separate.} Generation needs an accelerator; faithful
grading needs the scheduler and GNU coreutils. Generations travel as JSONL carrying the digest of
the task file, and the verifier refuses any whose digest does not match.

\section{Results}

Three seeds, $n=5$ samples per task, on an RTX 3060, graded inside the container. Every model was
loaded in 4 bit, which is the condition the tables carry unless a row says otherwise. The $\pm$ is
half the range across seeds. It is not a confidence interval, and no significance is claimed
anywhere in this paper.

\begin{table}[t]
\centering
\footnotesize
\setlength{\tabcolsep}{3pt}
\caption{T1, writing a job script from scratch. $\mathrm{pass}@1$ per level. A row is a
condition and not a model: the executor decides how \emph{functional} is judged and the
quantization decides how the model was loaded, and both are measurements rather than settings, so
each is published rather than one being chosen. Only the four other levels differ between the
\texttt{bash} and \texttt{sbatch} rows of one model. Safety is $1.000 \pm 0.000$ throughout and is
omitted.}
\label{tab:t1}
\begin{tabular}{lccccc}
\toprule
model, executor and loading & syntax & submit. & functional & res. fit & strict \\
\midrule
% generated by scripts/paper_data.py, do not edit
Qwen2.5-Coder 1.5B (4-bit) & $0.575 \pm 0.025$ & $0.842 \pm 0.013$ & $0.533 \pm 0.037$ & $0.442 \pm 0.013$ & $0.308 \pm 0.025$ \\
Qwen2.5-Coder 1.5B (fp16) & $0.750 \pm 0.000$ & $1.000 \pm 0.000$ & $0.533 \pm 0.025$ & $0.617 \pm 0.013$ & $0.408 \pm 0.025$ \\
Qwen3.5 9B (4-bit) & $1.000 \pm 0.000$ & $0.875 \pm 0.025$ & $0.650 \pm 0.025$ & $0.600 \pm 0.050$ & $0.450 \pm 0.025$ \\
Granite 4.1 3B (4-bit) & $1.000 \pm 0.000$ & $0.875 \pm 0.000$ & $0.842 \pm 0.037$ & $0.550 \pm 0.000$ & $0.425 \pm 0.000$ \\
Gemma 4 12B (4-bit) & $0.875 \pm 0.000$ & $0.867 \pm 0.013$ & $0.875 \pm 0.000$ & $0.917 \pm 0.050$ & $0.658 \pm 0.062$ \\
Qwen2.5-Coder 7B (4-bit) & $1.000 \pm 0.000$ & $0.792 \pm 0.025$ & $0.875 \pm 0.000$ & $1.000 \pm 0.000$ & $0.667 \pm 0.025$
 \\
\bottomrule
\end{tabular}
\end{table}

\begin{table}[t]
\centering
\footnotesize
\setlength{\tabcolsep}{3pt}
\caption{T2, diagnosing and repairing an induced fault. $\mathrm{pass}@1$ per level, one row per
condition as in Table~\ref{tab:t1}. Safety is $1.000 \pm 0.000$ throughout and is omitted.}
\label{tab:t2}
\begin{tabular}{lccccc}
\toprule
model, executor and loading & syntax & submit. & functional & res. fit & strict \\
\midrule
% generated by scripts/paper_data.py, do not edit
Qwen2.5-Coder 1.5B & $0.792 \pm 0.016$ & $0.886 \pm 0.005$ & $0.664 \pm 0.005$ & $0.412 \pm 0.014$ & $0.291 \pm 0.000$ \\
Granite 4.1 3B & $0.862 \pm 0.002$ & $1.000 \pm 0.000$ & $0.750 \pm 0.000$ & $0.753 \pm 0.009$ & $0.661 \pm 0.009$ \\
Qwen2.5-Coder 7B & $0.983 \pm 0.002$ & $0.977 \pm 0.000$ & $0.870 \pm 0.002$ & $0.965 \pm 0.007$ & $0.824 \pm 0.002$ \\
gemma-4-12B-it & $1.000 \pm 0.000$ & $0.876 \pm 0.002$ & $0.885 \pm 0.002$ & $0.938 \pm 0.002$ & $0.732 \pm 0.000$
 \\
\bottomrule
\end{tabular}
\end{table}

\begin{figure}[t]
\centering
\begin{tikzpicture}
\begin{axis}[
  anvilbar,
  xtick={0,1,2,3,4,5},
  xticklabels={syntax, submit., functional, res. fit, safety, strict},
  ylabel={$\mathrm{pass}@1$},
]
\addplot[fill=anvilA] table[x=idx, y=m0] {data/t1_slurm.dat};
\addplot[fill=anvilE] table[x=idx, y=m1] {data/t1_slurm.dat};
\addplot[fill=anvilB] table[x=idx, y=m2] {data/t1_slurm.dat};
\addplot[fill=anvilD] table[x=idx, y=m3] {data/t1_slurm.dat};
\addplot[fill=anvilC] table[x=idx, y=m4] {data/t1_slurm.dat};
\legend{Qwen2.5 1.5B, Qwen3.5 9B, Granite 3B, Gemma 12B, Qwen2.5 7B}
\end{axis}
\end{tikzpicture}
\caption{T1 per level, five models. Every level broadly rises with capability except
submittability, where the 3B, the 9B and the 12B lead and the 7B trails all four.}
\label{fig:t1}
\end{figure}

\subsection{Submittability is not ordered by size}

On T1 the level reads $0.875$ for a 3B model and for a 9B of a third generation, $0.867$ for a 12B,
$0.842$ for a 1.5B and $0.792$ for a 7B. Four models sit within eight points of each other at the
top and the 7B is alone at the bottom, while within the Qwen family the level falls as size rises
and every
other level improves, which Figure~\ref{fig:t1} shows as the one bar group running against the
others. Two families of three are above Qwen across a range from 3B to 12B, so the shortfall is not
one model's quirk.

T2 does not repeat that ordering, and it would be overclaiming to say otherwise: there every model
sits between $0.876$ and $1.000$, which leaves little room to separate them, and the family that
leads T1 is last.

The families also fail the level differently. One invents partition names, which a different cluster
might genuinely have; another invents an option SLURM does not have, which no cluster has. The
level measures how much a model volunteers that nothing asked for, and whose vocabulary it reaches
for when it does.

\subsection{Real submission was nearly redundant until a model wrote \texttt{srun}}

Grading the same generations twice inside the same image, once through a \texttt{bash} sandbox and
once through real \texttt{sbatch}, four models produced one changed verdict in $3120$ artifacts.
Everything real submission stopped had already failed another level, so the executor propagated a
verdict rather than producing one, and this section said so.

The fifth model ended that. Six artifacts of $3900$ now change verdict and $32$ samples fail under
the sandbox while passing under real submission, and five of the six share one cause:

\begin{center}
\small\texttt{srun: error: Unable to confirm allocation for job 1: Invalid job id specified}
\end{center}

The model writes \texttt{srun} inside the batch script, which is how a job step is launched on a
real cluster. The sandbox has no allocation for it to attach to, so the step exits non-zero and the
sample is recorded as failing. The sandbox produces a false negative on every script that launches
a job step, and did so silently while no model wrote one.

Both artifacts the two executors disagree about run the same direction, and it is not the one a
sandbox is usually suspected of: neither is the scheduler catching something the sandbox missed,
both are the sandbox rejecting what the scheduler accepts, for want of an allocation or a core
binding. The deflationary claim was true of the four models measured and not of the method, which
is the more useful way to have been wrong.

What it was really a claim about is the task set. Every task in the two sets above states its
resource requirements in its prompt, so a static check can reach them and the executor has little
left to decide. A second set states no memory minimum: the payload is written by the model, what it
needs is a property of that payload, and no level that reads the text can know it. Graded the same
way, five models and three seeds and $900$ sample comparisons, the same executor comparison reads
\textbf{$288$ artifacts of $900$}, with $356$ of the $394$ stopped artifacts returning
\texttt{OUT\_OF\_MEMORY}. From $0.15\%$ to $32\%$, on the same harness and the same models.
That figure is one grading. Repeated under five successive verifiers it read $290$, $289$, $288$,
$289$ and $288$, moving by the single artifact discussed in the previous section; the sample it
turns on is a repair of a memory under-request sitting against the allocation it asked for, which
is what this set exists to exercise.

Table~\ref{tab:exec} gives the from-scratch half, one row per arm. Qwen2.5-Coder 7B is the case:
the sandbox promotes $29$ of its $30$ artifacts and the scheduler kills every one, so
\emph{functional} falls from $0.967$ to $0.000$ between its two rows. It is the model with $1.000$ resource fit on the main T1 set, strong
wherever the requirement is written down and at floor wherever the requirement is a property of what
it wrote itself. Gemma 4 12B is the only model that solves the set, $1.000$ under both arms.

\begin{table}[t]
\centering
\footnotesize
\setlength{\tabcolsep}{3pt}
\caption{The execution-sensitive set from scratch, each model graded twice, $\mathrm{pass}@1$ per
level. Two tasks, three seeds, $n=5$. The set states no memory minimum, so what a script needs is a
property of the payload the model wrote and only running the job can decide it. Safety is
$1.000 \pm 0.000$ throughout and is omitted.}
\label{tab:exec}
\begin{tabular}{lccccc}
\toprule
model, executor and loading & syntax & submit. & functional & res. fit & strict \\
\midrule
% generated by scripts/paper_data.py, do not edit
Qwen2.5-Coder 1.5B & $1.000 \pm 0.000$ & $0.500 \pm 0.200$ & $0.167 \pm 0.100$ & $0.667 \pm 0.050$ & $0.133 \pm 0.050$ \\
Qwen3.5 9B & $0.833 \pm 0.100$ & $0.833 \pm 0.100$ & $0.467 \pm 0.050$ & $1.000 \pm 0.000$ & $0.467 \pm 0.050$ \\
Granite 4.1 3B & $1.000 \pm 0.000$ & $1.000 \pm 0.000$ & $0.467 \pm 0.200$ & $0.133 \pm 0.050$ & $0.067 \pm 0.100$ \\
Gemma 4 12B & $1.000 \pm 0.000$ & $1.000 \pm 0.000$ & $1.000 \pm 0.000$ & $1.000 \pm 0.000$ & $1.000 \pm 0.000$ \\
Qwen2.5-Coder 7B & $1.000 \pm 0.000$ & $0.500 \pm 0.000$ & $0.000 \pm 0.000$ & $0.500 \pm 0.000$ & $0.000 \pm 0.000$
 \\
\bottomrule
\end{tabular}
\end{table}

Qwen3.5 9B moves the other way for the reason above, \texttt{srun} in the payload. On the repair
half every model loses at least a third of its strict score under enforcement while the ordering is
preserved, so this is a level of difficulty the sandbox cannot see rather than a re-ranking. The set
is small, two tasks and ten repairs, and its numbers are not comparable with the $3900$-sample
tables; what it establishes is not a score but that the executor's value is a property of what is
being asked.

\subsection{Retrieval costs, and the sign depends on prior knowledge}

\begin{figure}[t]
\centering
\begin{tikzpicture}
\begin{axis}[
  anvilbar,
  bar width=11pt,
  height=5.6cm,
  xtick={0,1,2,3},
  xticklabels={no context, placement doc, time/mem doc, off-topic},
  ylabel={resource fit $\mathrm{pass}@1$},
]
\addplot[fill=anvilA] table[x=idx, y=small] {data/retrieval.dat};
\addplot[fill=anvilC] table[x=idx, y=large] {data/retrieval.dat};
\legend{Qwen2.5-Coder 1.5B, Qwen2.5-Coder 7B}
\end{axis}
\end{tikzpicture}
\caption{Resource fit under four retrieval conditions. The small model pays a toll for any attached
text, including a length-matched passage on an unrelated subject; the large model pays none, and
loses ground only to the document about the two directives this level checks.}
\label{fig:retrieval}
\end{figure}

Seven conditions across two sizes, four of them shared and plotted in
Figure~\ref{fig:retrieval}. Any attached text costs the 1.5B about $28$ points of resource fit,
whether it is SLURM documentation or a passage about coastal tides cut to the same length. The 7B
pays no such toll: two unrelated off-topic controls leave it on its zero-shot values to the digit.
The single corpus document addressing \verb|--time| and \verb|--mem| is worth $+21$ points to the
model that does not know those directives have no default, and $-6$ to the one that does.

Capacity decides whether a model can ignore what does not concern it, and relevance decides the
sign of what it cannot ignore.

\subsection{A toolchain difference no model reaches}

GNU coreutils 9.4 and \texttt{uutils} 0.8.0 agree on $91$ of $101$ probed invocations. Three of the
disagreements are behavioural and need no misconfigured locale: character counting and column
expansion on a non-ASCII payload, and SI number formatting. A task built to sit in that corner is
judged differently by the two implementations, which a guard asserts.

Asked to solve it, three models produced $45$ samples: none pinned a locale, and none diverged.
Each fails on both implementations for a reason that precedes the difference, one of them by
emitting escape sequences as literal text and so removing the non-ASCII content the question was
about.

\section{What went wrong while measuring}

Reported because the corrections are load-bearing.

A published multi-seed table had been graded against the experiment machine's own scheduler, which
has one node, no GPUs, and rejects three of the eight canonical solutions. Submittability and
strict were wrong by up to $36$ points, while syntax, functional and resource fit were unchanged to
the digit, which is the signature that identified the fault. The topology canary exists because of
it, and fired on that same machine the first time a model was run after it landed.

The same defect was later found in the retrieval ablation, whose two scheduler-dependent columns
were withdrawn and remeasured. Correcting them inverted an arm ordering: a single-seed pilot had
been right, and the larger run that appeared to correct it had been wrong. The cause was the
grading environment, not the seeds.

Three mechanisms proposed for the retrieval effect were each refuted by the experiment run to test
them. The fourth was found only by counting per sample instead of per problem.

The resource fit check compared the requested walltime against the one a task names from above
only, so a script asking fifteen seconds where the prompt asked fifteen minutes was scored correct.
$123$ of $2421$ passing artifacts were requests of that kind. Regrading the saved generations under
a floor moved five of ten published cells, left the other five identical to the digit, and dropped
one model below another on strict: a one-sided comparison does not inflate a table evenly, it
inflates whichever model has the habit the check cannot see. Reports and leaderboard entries now
carry a digest of the verifier beside the digest of the task file, since a changed check
invalidates a comparison exactly as thoroughly as a changed task and left no trace of having done
so.

That gap surfaced while models were being screened for something else, so the remaining constraints
were then audited on purpose. Memory had the same shape, compared from below only, and $30$ of
$2298$ passes requested more than the task named, all of them repairs in which the model rewrote a
directive the induced fault had not touched. Tightening it moved one cell, by the amount the audit
had predicted to three decimals before the change was made. Requested GPU counts were exact in all
$343$ cases, so the measurement asked for no correction; the check was made two-sided anyway, on
the grounds that a hole nobody has fallen into is still a hole, and regrading confirmed that it
moved nothing. All six constraints carrying a value now demand equality.

\section{Threats to validity}

Eight T1 tasks, which makes each one worth $0.125$ of every T1 figure. Three seeds, so the ranges
are spread rather than intervals. Three model families at five sizes, from 1.5B to 12B, quantized to
4 bit except where a row says otherwise, and that exception is not decoration: the 1.5B loaded in
fp16 gains $17.5$ points of syntax, $15.8$ of submittability and $17.5$ of resource fit while
\emph{functional} does not move at all, so the loading is part of what a row measures rather than a
setting behind it. Only the 1.5B fits in fp16 on this card, so how far that generalises is unknown.
One of the five models reasons by default and was measured with that disabled, since under this
token budget it would otherwise be cut off before producing a script. One reference topology,
declared by this work rather than borrowed from a centre in production. All prompts written by one
author, in one language and one style. No result here has
been replicated independently.

One figure per cell, which is less of a guarantee than it looks. Verifying one set of generations
twice, under one verifier and one container image, returned two different values for
\emph{functional}: it is the only level that executes, so it is the only one that can disagree with
itself, and the digest of the verifier bounds which rules produced a verdict without bounding the
verdict. The instance behind that observation was a race in the harness and is fixed, but the
property is general and the fix is not a proof that it was the only one. Under real submission a
job sitting against the allocation it requested is killed in one grading and not the next, five
gradings running, which is arguably the measurement working rather than failing. Every figure here
comes from a single grading, and a difference of a few thousandths between two rows should be read
with that in mind.

\section{Availability}

Code, tasks, canonical solutions, the digest manifest identifying the dataset version, the
generated leaderboard and every figure above are at
\url{https://github.com/antoniorotundo2/anvil} under MIT. The command \texttt{anvil check} applies
the verifier to a single artifact, with or without a site policy, for readers who want the
instrument rather than the measurements.

\bibliographystyle{plainnat}
\bibliography{anvil}

\end{document}
